\section{Implementation}

The target environment is the module-provided Ubuntu VM with Python 3, Mininet, Ryu, and Open vSwitch. Exact installed versions and VM kernel details are to be recorded from the validated run and inserted here: [Ubuntu/kernel: pending]; [Python: pending]; [Mininet: pending]; [Ryu: pending]; [Open vSwitch: pending]. The topology configures OpenFlow 1.3 and uses the controller address and port shown in the coursework figure.

The code is divided by responsibility. \texttt{topology.py} builds and starts the Mininet network. \texttt{learning\_switch.py} is a Task 1-only controller without the guest restriction. \texttt{controller.py} handles switch events, per-switch MAC learning, policy enforcement, five-tuple detection, CSV alert logging, and timed mitigation. \texttt{flow\_detection.py} is a small standard-library rolling-window component that can be tested without Ryu. \texttt{traffic\_generator.py} produces ICMP echo attempts or TCP connection attempts at a chosen rate and duration. TCP tests use a Python HTTP service on h7. \texttt{monitoring.py} polls flow and port counters and produces CSV and labeled SVG graphs without requiring an extra package; it also produces PNG graphs if Matplotlib is already available. \texttt{experiment.py --live} starts the controller, topology, traffic, monitoring, and five-run evaluation in the VM, then writes results and a transcript to unique output paths.

The detection window and threshold were selected to distinguish the demonstrated rates: one new connection per second remains below the threshold, whereas an eight-per-second generator should exceed it within a five-second window. A 15-second hard timeout is long enough to verify the block while permitting automatic recovery during a short demonstration. These settings require validation against the actual VM; the threshold may need adjustment if connection attempts are delayed or coalesced by the environment. The implementation handles the Task 4 mitigation and guest rule at s3 because h6, h7, and h8 attach to that switch.

The traffic generator accepts a mode, destination address, rate, duration, and TCP port. ICMP mode invokes one-packet ping attempts at the requested cadence; TCP mode repeatedly opens and closes client connections. The high-rate experiment is therefore a connection-creation workload, not a sustained bulk-data transfer. That distinction matters when interpreting counters: short connections increase the number of flow entries and packet-in events, while the total byte count can remain modest. The web server on h7 provides a listening endpoint for connection establishment and can be started without installing a separate service package.

\texttt{monitoring.py} runs outside the Mininet host namespaces and queries each named OVS bridge. Flow output provides packet and byte counters, durations, and match entries; port output provides receive and transmit packet counts. Active flows exclude the priority-zero table-miss entry. The monitoring CSV has one flow row and one port row per switch and sample time. The active-flow graph plots the count for each switch, while the traffic-volume graph plots cumulative bytes. Both graphs use sample number on the horizontal axis so that irregular VM scheduling does not imply a more precise clock than the data supports. The SVG writer uses only the standard library; PNG output is supplemental when Matplotlib is preinstalled in the module VM.

The integrated runner records unique output files so that successive experiments do not silently replace earlier results. Its transcript records host addresses, bridge protocols, initial flows, policy checks, each performance run, and output locations. The Ryu log is retained separately to make the alert timestamp and switch connection messages auditable. For a formal report, the evidence set should keep the CSVs and graphs alongside readable terminal screenshots rather than relying on the runner's summary alone.

One design issue is the relationship between flow visibility and forwarding. A broad TCP flow rule would allow later connections to bypass packet-in and make new-connection counting incomplete. The controller therefore matches TCP source and destination ports as well as IP addresses and protocol. This yields a packet-in for a new five-tuple while letting later packets in the same connection use the installed rule. The trade-off is more short-lived flow entries, controlled here by the five-second idle timeout.
